\begin{textsample}{POS Dim 1 – sp – Score 39.00 – sp/sp\_em\_30.txt}  \label{ex:f1_pos_116}
FÍSICA O componente curricular de Física, na área de \textbf{Ciências} da Natureza e suas Tecnologias, propõe ir além do aprendizado de seus conceitos \textbf{teóricos}.
A aprendizagem de Física no Ensino Médio, qualitativamente distinta da realizada no Ensino Fundamental, deve aprofundar os conceitos, estruturando leis e princípios reconhecidos e estabelecidos para a incorporação de novas aprendizagens, contextualizando os objetos de conhecimento existentes, por meio de demonstrações e experimentos científicos que simulem situações cotidianas e promovam a autonomia e o protagonismo do estudante utilizando, para isso, metodologias ativas de ensino e investigação científica.
A aprendizagem no componente Física indica a \textbf{compreensão} do mundo físico e natural e o desenvolvimento de estratégias que apliquem os conhecimentos para investigar e propor ações de intervenção no mundo \textbf{contemporâneo}.
O objetivo é \textbf{aproximar} o estudante da investigação científica e tecnológica para o enfrentamento dos desafios cotidianos, na direção de uma educação integral, explorando diferentes modos de pensar e falar sobre a cultura científica, presente na diversidade sociocultural de cada estudante, promovendo a aprendizagem em diversos contextos históricos, sociais e étnicos.
Os temas \textbf{transversais} podem ser utilizados para o trabalho interdisciplinar entre os componentes curriculares da área.
É possível realizar ações relacionadas a temas como : Ciência e Tecnologia, Meio Ambiente, Economia, Saúde e Cidadania.
De acordo com as Orientações Educacionais Complementares aos \textbf{Parâmetros} Curriculares Nacionais ( PCN+ Ensino Médio, 2006 ), cada área do conhecimento tem sua importância no ensino das \textbf{ciências} físicas : as Linguagens e a Matemática permitem representar e comunicar por meio de textos e demonstrações matemáticas ; a área de Ciências Humanas constrói a evolução dos conceitos da Física e contextualiza as inovações tecnológicas da sociedade, discutindo o uso ético e \textbf{socioambiental} dessas inovações ; a área de \textbf{Ciências} da Natureza, a qual da Física faz parte, incentiva a investigação e \textbf{compreensão} do mundo \textbf{atual}, juntamente com os fenômenos químicos e biológicos.
Esse engajamento entre as áreas permite ao estudante do Ensino Médio a construção de novos saberes e o desenvolvimento de \textbf{competências} e habilidades para solucionar \textbf{problemas} do mundo real.
É importante ressaltar que os temas da Física estão contemplados nos objetos de conhecimento presentes no \textbf{organizador} da Base Comum, porém o desenvolvimento desses objetos será feito de forma a apresentar e experienciar os conceitos físicos relacionados aos fenômenos da natureza e, posteriormente, aprofundar os objetos de conhecimento nos itinerários formativos envolvendo Ciências da Natureza e suas Tecnologias.
Para organizar os objetos de conhecimento de Física para o Currículo do Ensino Médio, foi realizado um estudo sobre todas as habilidades existentes no Currículo Paulista da etapa do Ensino Fundamental relacionadas à Física.
No Ensino Fundamental, especificamente no 7º ano, é realizada a discussão e a investigação de máquinas \textbf{simples} na realização de \textbf{tarefas} mecânicas, a fim de propor soluções para a sociedade.
Conceitos de mecânica são estudados utilizando a Astronomia como objeto de conhecimento para o desenvolvimento de habilidades.
Os objetos de conhecimento relacionados à Astronomia se encontram em todos os ciclos de aprendizagem do Ensino Fundamental.
Nos anos iniciais ( alfabetização - 1º ao 3º ano ), as habilidades de observar e registrar aparecem com \textbf{certa} frequência, visto que o estudante nesta etapa, inicia seu processo de alfabetização nas hipóteses de escrita existentes, utilizando o objeto de conhecimento relacionado à observação do céu e aos efeitos da radiação solar em diferentes superfícies.
Após a observação e registro, o estudante inicia a identificação e \textbf{descrição} dos corpos celestes nos ciclos diários, além de reconhecer os avanços tecnológicos existentes.
Do 4º ao 6º ano o estudante compara e analisa os pontos cardeais e inicia a explicação dos movimentos dos corpos celestes, associando -os à marcação de tempo, fases da Lua e identificação de constelações e movimentos aparentes de estrelas, reconhecendo e explicando os movimentos terrestres.
Já do 7º ao 9º ano, o estudante é incentivado a construir modelos do Sistema Solar e das fases da Lua para representá -los.
Ao \textbf{final} do Ensino Fundamental, além de aprofundar os conceitos de corpos celestes, é possível investigar os avanços tecnológicos discutindo a exploração espacial e a possível sobrevivência fora da Terra, por meio de pesquisas sobre as características de outros planetas e sobre viagens planetárias e interestelares.
Os objetos de conhecimento relacionados à Termologia são desenvolvidos no 4º, 7º e 9º anos, iniciando pela investigação das transformações dos materiais de acordo com as mudanças térmicas, \textbf{verificadas} por meio de experimentos científicos para concluir a reversibilidade na alteração de estado físico, reconhecendo tais mudanças em fenômenos do cotidiano.
Na \textbf{sequência}, o estudante desenvolve habilidades para diferenciar conceitos sobre temperatura, calor e sensação térmica, permitindo construir soluções e explicar o funcionamento de equipamentos térmicos, utilizando conceitos de propagação de calor \textbf{tanto} para justificar a utilização de materiais condutores e isolantes, quanto para identificar e analisar o equilíbrio termodinâmico para a manutenção da vida na Terra.
O objetivo é \textbf{argumentar} em discussões sobre os avanços tecnológicos e as consequências \textbf{socioambientais} na utilização dos conceitos e produção de máquinas e materiais.
A partir daí, é possível ao estudante entender a estrutura da matéria de acordo com seu estado físico, explicando e representando tais transformações com base na constituição submicroscópica.
Em Luz e Som, as habilidades são desenvolvidas durante o 3º, 5º, 6º e 9º anos.
A \textbf{abordagem} inicial desses objetos de conhecimento desenvolve a identificação, \textbf{descrição} e experimentação de conceitos iniciais sobre a luz e o som, além da produção de diferentes sons e do reconhecimento de condições prejudiciais à saúde humana.
Após a \textbf{abordagem} inicial, o estudante é incentivado a planejar e executar experimentos científicos relacionados à luz e ao som e, também, a analisar e explicar como ocorre a transmissão e \textbf{recepção} da imagem e do som ( radiações eletromagnéticas ) e compreender os avanços tecnológicos envolvendo as radiações em nossa vida, com a utilização de aparelhos eletrônicos pessoais, com os aparelhos de \textbf{exames} laboratoriais e no tratamento de doenças.
A temática Eletricidade é estudada durante o 8º ano, construindo -se circuitos elétricos para aprender sobre seus componentes e utilização em circuitos residenciais, classificando -se os equipamentos de acordo com as transformações de energia, calculando -se o consumo de eletrodomésticos a partir da potência e do tempo de uso do equipamento, a fim de comparar e avaliar o impacto na economia de energia elétrica para o consumo sustentável.
Após o estudo e a \textbf{análise} dos equipamentos elétricos e eletrônicos, o estudante identifica, classifica e avalia a produção de energia elétrica de fontes renováveis e não renováveis para propor soluções para o uso consciente da produção e do consumo dessa energia, utilizando princípios de sustentabilidade e acessibilidade \textbf{econômica}.
O estudante tem a possibilidade de discutir, investigar e avaliar os avanços tecnológicos de informação e automação por meio de transmissão e acesso de dados para \textbf{argumentar}, de forma \textbf{crítica} e \textbf{reflexiva}, sobre a natureza das informações no mundo do trabalho e na vida pessoal.
No Ensino Médio, o estudante terá a oportunidade de ampliar e aprofundar seus conhecimentos através dos objetos de conhecimentos contemplados nas unidades temáticas : Matéria e Energia ; Vida, Terra e Cosmos ; \textbf{Tecnologia} e Linguagem Científica.
O desenvolvimento da unidade temática Matéria e Energia, em Física, se dará por meio do estudo das relações e interações entre matéria e energia que estão presentes em fenômenos naturais e nos processos tecnológicos.
Isso permitirá ao estudante avaliar os impactos dessas interações no mundo cotidiano através da aplicação de modelos e propostas de intervenção em diversos contextos, sendo capaz de desenvolver \textbf{competências} relacionadas à \textbf{compreensão} das estruturas microscópicas da matéria.
Na unidade temática Vida, Terra e Cosmos, a Física, propõe estudos científicos que fundamentam as \textbf{teorias} e leis sobre a origem da vida, do planeta, do universo e das interações gravitacionais.
A relação com a área de Matemática e suas Tecnologias, permitirá ao estudante elaborar explicações e previsões de cálculo para os movimentos dos corpos na Terra e no universo.
A proposta da Física para o desenvolvimento da unidade temática \textbf{Tecnologia} e Linguagem Científica é que o estudante \textbf{mobilize} e articule os conhecimentos científicos acerca do uso das \textbf{tecnologias} como meio de comunicação e para a resolução de \textbf{problemas} sociais, \textbf{econômicos} e sustentáveis, e avalie os avanços tecnológicos que contribuem para uma melhor qualidade de vida.
Dessa forma, a Física deve ter como premissa permitir que o estudante compreenda o mundo \textbf{moderno} e \textbf{contemporâneo}, seus desafios e as possibilidades que a sociedade do conhecimento oferece para representar esse mundo.
O desenvolvimento de \textbf{competências} e habilidades deve ser expresso em torno de assuntos e \textbf{problemas} reais, que \textbf{exigem} aprendizagem de leis, conceitos e objetos de conhecimento construídos através de um processo cuidadoso de identificação das relações do conhecimento científico.
Atualmente, na sociedade, ouve -se música digitalizada, manuseiam -se computadores que \textbf{operam} com semicondutores, dados são transmitidos por meio de ondas eletromagnéticas nos smartphones, a medicina dispõe de aparelhos de raios-x e ressonância magnética, e o laser revolucionou as técnicas médicas e odontológicas.
Os aceleradores de partículas de altíssimas energias estabelecem uma abertura para o entendimento das estruturas microscópicas do espaço-tempo, permitindo, assim, decifrar as leis fundamentais que regem o comportamento da matéria, possibilitando o entendimento da origem do Universo e \textbf{desvendar} fenômenos físicos.
A Física se torna um importante componente no desenvolvimento das \textbf{tecnologias} \textbf{digitais} de informação e comunicação ( TDICs ), do \textbf{letramento} \textbf{digital} e pensamento computacional, eixos estruturantes do componente curricular de Tecnologia e Inovação.
Nesse contexto, a \textbf{tecnologia} colabora para o \textbf{diagnóstico}, a formação integral e a avaliação do estudante.
Portanto, a aprendizagem no componente de Física indica a \textbf{compreensão} e a utilização dos conhecimentos científicos, para analisar e explicar o funcionamento do mundo e para propor ações de intervenção de modo sustentável, com o auxílio das \textbf{tecnologias}.
Assim, a Física se mostra um importante componente da área para explicar o mundo tecnológico \textbf{atual}.

% matched lemmas: abordagem, análise, aproximar, argumentar, atual, certo, ciência, competência, compreensão, contemporâneo, crítico, descrição, desvendar, diagnóstico, digital, econômico, exame, exigir, final, letramento, mobilizar, moderno, operar, organizador, parâmetro, problema, recepção, reflexivo, sequência, simples, socioambiental, tanto, tarefa, tecnologia, teoria, teórico, transversal, verificar
\end{textsample}
